% !TeX root = ./../../main.tex
% Schemat blokowy obwodu realizującego szereg dwuskładnikowy (1/2)(kappa_{1,0} + kappa_{2,0}).
% Rysunek własny (quantikz2); struktura zgodna z obwodem z pracy magisterskiej autora
% (artykuły/magisterka_werner_final/graf/Eksperymenty/wiele_komp/uklad.png).
% Obwód jest składany w pudełku i skalowany do szerokości otoczenia (\linewidth).
\ifdefined\tbqsampbox\else\newsavebox{\tbqsampbox}\fi
\begin{lrbox}{\tbqsampbox}
\begin{quantikz}[row sep=0.2cm, column sep=0.2cm, font=\small]
\lstick{$b\colon\ket{0}$} & \gate{\mathrm{H}} & \gate[2]{\compphase^{(b)}_{2,0}} & & \gate[3]{\bigl(\compphase^{(ab)}_{3,0}\bigr)^{\dagger}} & \ctrl{2} & & \\
\lstick{$\ket{D}$} & \qwbundle{m} & & \gate[2]{\compphase^{(a)}_{1,0}} & & & & \\
\lstick{$a\colon\ket{0}$} & \gate{\mathrm{H}} & & & & \targ{} & \gate{\mathrm{H}} & \meter{}
\end{quantikz}
\end{lrbox}
\resizebox{\linewidth}{!}{\usebox{\tbqsampbox}}
